\section*{Chapitre \ref{ch:dv:fx-derivatives} --- Dérivés de change}

\begin{solution}{exo:dv:fx-derivatives:1}
$e^{-r_fT}S=e^{-r_dT}F$, donc $e^{-r_fT}S\Phi(d_1)-e^{-r_dT}K\Phi(d_2)=e^{-r_dT}\bigl(F\Phi(d_1)-K\Phi(d_2)\bigr)$, la formule de Black appliquée au forward avec
le facteur d'actualisation de la devise de cotation~; $d_1$ et $d_2$ ne dépendent de $S$ qu'à travers $F$.
\end{solution}

\begin{solution}{exo:dv:fx-derivatives:2}
La crainte du marché est une chute de la devise émergente, une hausse du dollar face à elle, et de telles chutes s'accompagnent d'une volatilité croissante~: le taux et sa
variance évoluent ensemble, $\rho>0$. Les calls sur le dollar (la protection) sont alors cotés à des volatilités plus élevées que les puts, et le risk reversal est
positif~: $+1.89$ points à 25 delta.
\end{solution}

\begin{solution}{exo:dv:fx-derivatives:3}
Une position vendeuse d'options à effet de levier~: à chaque fixing au-dessus du strike, le client vend deux fois le notionnel à un taux inférieur au marché, et la cible plafonne
les gains~; l'amélioration correspond à la prime de cette position, reçue sous forme d'un meilleur taux.
\end{solution}

\begin{solution}{exo:dv:fx-derivatives:4}
Les poids correspondent exactement au \omterm{def:dv:greeks-and-the-hedging-pnl:greeks}{vega}, au \omterm{def:dv:greeks-and-the-hedging-pnl:greeks}{vanna} et au \omterm{def:dv:greeks-and-the-hedging-pnl:greeks}{volga} des piliers~: à un pilier, le poids vaut un sur lui-même et zéro ailleurs, et le prix est celui
du marché. Au-delà des piliers, la méthode suppose que le coût du smile est porté par ces trois risques à la volatilité à la monnaie~; les ailes à 10 delta présentent une
courbure que les cotations à 25 delta ne déterminent pas, si bien que l'extrapolation la manque (un écart au call à 10 delta ici).
\end{solution}

\begin{solution}{exo:dv:fx-derivatives:5}
En \omterm{def:dv:local-volatility:model}{volatilité locale}, la volatilité est une fonction du spot, et l'aile d'achat du marché est chère~; ainsi, lorsque le taux monte vers 7.20, sa volatilité
augmente avec lui et il atteint la barrière plus facilement. Dans Heston, une partie de la volatilité est portée par un processus de variance qui n'est pas lié au spot, si bien que
le chemin vers la barrière est en moyenne moins volatil.
\end{solution}

\begin{solution}{exo:dv:fx-derivatives:6}
Avec $\lambda=0$, la variance suit $dv=\kappa(\bar v-v)dt$, une trajectoire déterministe $v(t)=\bar v+(v_0-\bar v)e^{-\kappa t}$~; alors
$\E[v_t\mid S_t=S]=v(t)$ et $L(t,S)^2v(t)=\sigma_{\mathrm{loc}}(t,S)^2$.
\end{solution}

\begin{solution}{exo:dv:fx-derivatives:7}
Le strike à coût nul vaut 7.325 en \omterm{def:dv:local-volatility:model}{volatilité locale}, 7.312 en \omterm{def:dv:local-volatility:model}{volatilité locale} stochastique et 7.305 dans Heston~; le nombre espéré de fixings est de 1.64,
1.73 et 1.78. Le modèle dont le chemin vers le haut est le plus volatil (la \omterm{def:dv:local-volatility:model}{volatilité locale}) valorise plus cher les calls vendus à effet de levier du client, si bien que le dealer
peut offrir un meilleur strike.
\end{solution}

\begin{solution}{exo:dv:fx-derivatives:8}
Les 98~\% se rachètent avec un gain faible et plafonné (0.30 par dollar au plus)~; les 2~\% restants courent avec deux fois le notionnel contre le client et portent toute la
valeur du produit~: la perte espérée sur ces trajectoires est assez grande pour que l'ensemble vaille zéro. La probabilité de rachat ne dit rien de l'ampleur des pertes.
\end{solution}

\begin{solution}{pb:dv:fx-derivatives:1}
\textbf{1.} 5.32~\%~; $+1.89$ et $+0.46$ points.
\textbf{2.} Que la hausse du dollar face à la devise, c'est-à-dire une dévaluation, soit le scénario redouté, valorisé à une volatilité plus élevée.
\textbf{3.} 6.861~: le taux du dollar (5~\%) est supérieur à celui de la devise (3~\%), si bien que le forward est inférieur au spot du montant du carry.
\textbf{4.} À un dixième de point près à 10 et 25 delta et à la monnaie.
\textbf{5.} 0.79.
\textbf{6.} 7.312, soit 6.6~\% au-dessus du forward.
\textbf{7.} 98.1~\% se rachètent, après 1.73 fixings en moyenne.
\textbf{8.} Un gain de 0.30 million en un ou deux fixings sur la plupart des trajectoires (95~\% gagnent au moins 126\,000)~; la valeur nulle provient des rares trajectoires sur lesquelles le
taux monte et reste au-dessus du strike, ce qui coûte chaque mois deux fois le notionnel.
\textbf{9.} 7.325 en \omterm{def:dv:local-volatility:model}{volatilité locale}, 7.305 dans Heston.
\textbf{10.} Une cible plus élevée permet au client de gagner davantage avant le rachat, si bien que le produit vit plus longtemps et expose moins le dealer de marché du côté
à effet de levier, pour un strike donné~; pour rester à coût nul, le strike doit être moins généreux.
\textbf{11.} 7.12 millions par million de notionnel mensuel, dont 1.01 million au premier trimestre.
\textbf{12.} 0.99 million.
\textbf{13.} 11.7 fixings~: après une hausse de 10~\%, le produit ne se rachète presque jamais.
\textbf{14.} Son delta change brutalement~: près de la cible, les gains restants sont faibles et le produit est sur le point de disparaître, si bien que la couverture des futurs
forwards à effet de levier doit être dénouée~; le delta peut changer de signe en quelques pips.
\textbf{15.} Une dépréciation par paliers de la devise émergente dans un régime managé est le scénario sur lequel repose la valeur du TARF~: une chute de 2.8~\% en deux jours après un
changement de règle de fixing relève de ce qui transforme un produit censé se racheter aussitôt en un produit qui va jusqu'au bout avec deux fois le notionnel.
\textbf{16.} La distribution des issues, et pas seulement la probabilité de rachat~: la perte dans les scénarios qui comptent (un mouvement de 10~\% coûte 7.1 millions par
million), ainsi que l'effet de levier.
\textbf{17.} La dynamique du smile et les sauts~: le strike à coût nul se déplace de 0.02 (200 pips) d'un modèle à l'autre, mais une dévaluation par paliers échappe aux trois.
\textbf{18.} Non~: sans revenu naturel en dollars, les ventes à effet de levier sont spéculatives et non une couverture~; un régulateur et les règles de convenance de la banque y feraient objection.
\textbf{19.} 1.73 fixings en moyenne avant le rachat~; une perte espérée de 7.12 millions par million de notionnel mensuel si le taux monte de 10~\% au cours du premier
trimestre.
\textbf{20.} Un strip d'options à effet de levier sur un taux qui évolue contre le client, rémunéré par un meilleur taux et des gains plafonnés.
\end{solution}

\begin{solution}{iq:dv:fx-derivatives:1}
À la monnaie~: la volatilité du straddle delta neutre. Risk reversal~: la volatilité du call à 25 delta moins celle du put (le skew). Butterfly~: la moyenne des ailes
par rapport à la monnaie (la courbure), selon la convention de la paire (smile ou strangle de courtier). Convertir avec la convention delta de la paire en trois
points strike--volatilité, puis interpoler (vanna--volga, SABR par échéance, ou un smile paramétrique).
\iqlookfor{les trois cotations et les conventions.}
\end{solution}

\begin{solution}{iq:dv:fx-derivatives:2}
Résoudre les poids des trois piliers qui égalent le \omterm{def:dv:greeks-and-the-hedging-pnl:greeks}{vega}, le \omterm{def:dv:greeks-and-the-hedging-pnl:greeks}{vanna} et le \omterm{def:dv:greeks-and-the-hedging-pnl:greeks}{volga} de l'option à la volatilité à la monnaie~; ajouter au prix de Black les coûts pondérés du marché moins Black
des piliers. Elle suppose que le coût du smile est entièrement capturé par ces trois risques, valorisés par les instruments cotés, une «~volatilité plate mais
stochastique~».
\iqlookfor{les poids et l'hypothèse.}
\end{solution}

\begin{solution}{iq:dv:fx-derivatives:3}
$dS/S=\dots+L(t,S)\sqrt v\,dW$ avec une variance stochastique~: la \omterm{def:dv:fx-derivatives:leverage}{fonction de levier} permet de caler chaque vanille, le processus de variance fournit la dynamique.
Calibrer $L(t,S)^2=\sigma_{\mathrm{loc}}^2/\E[v\mid S]$ par la méthode particulaire~: simuler, estimer l'espérance conditionnelle à partir des particules à chaque pas,
mettre à jour $L$~; choisir le mélange de la volatilité de la variance à partir des prix exotiques.
\iqlookfor{la formule de projection et la méthode particulaire.}
\end{solution}

\begin{solution}{iq:dv:fx-derivatives:4}
La différence tient à la dynamique, et non au smile~: calibrer le paramètre dynamique du modèle (le poids de mélange) sur les prix de touche du marché là où ils se traitent~;
à défaut, valoriser avec le modèle dont la dynamique est soutenue par le comportement historique du smile, et conserver une réserve pour l'écart entre les
modèles plausibles.
\iqlookfor{calibrer la dynamique sur les exotiques, et une réserve.}
\end{solution}

\begin{solution}{iq:dv:fx-derivatives:5}
Le client vend une devise à terme à un meilleur taux tant qu'il gagne, avec des gains plafonnés par une cible, et deux fois le montant à un taux moins favorable tant qu'il perd, sans
plafond. Pour le client~: de fortes pertes dans un mouvement soutenu contre lui. Pour la banque~: la sensibilité de la couverture à la dynamique du smile et aux sauts, et le
risque de crédit et de convenance du client à mesure que les pertes s'accumulent.
\iqlookfor{l'asymétrie, et les risques des deux côtés.}
\end{solution}

\begin{solution}{iq:dv:fx-derivatives:6}
Ajouter une composante de saut (ou de changement de régime) à la valorisation et aux scénarios~: des barrières peuvent être franchies sans transaction à proximité, et un TARF peut passer du rachat
à une exposition au double du notionnel du jour au lendemain. Couvrir avec des options qui paient lors du saut, limiter la concentration aux niveaux de barrière proches du bord d'une bande, et constituer une réserve pour le
saut au niveau du livre.
\iqlookfor{sauts dans la valorisation, couvertures de saut et réserves.}
\end{solution}
